\begin{thebibliography}{99}
\bibitem{ABL}
E. Assing, V. Blomer and J. Li,
\emph{Uniform Titchmarsh divisor problems},
\href{https://arxiv.org/abs/2005.13915v1}{arXiv:2005.13915v1} (2020).

\bibitem{ChanSupplement}
S. Chan,
\emph{Theorem $(1+1.75)$ on the Goldbach Conjecture --- Numerical Supplement},
version 20261106, Zenodo.
\href{https://doi.org/10.5281/zenodo.23755453}{doi:10.5281/zenodo.23755453}.

\bibitem{Drappeau}
S. Drappeau,
\emph{Sums of Kloosterman sums in arithmetic progressions, and the error
term in the dispersion method}, Proc. Lond. Math. Soc. (3)
\textbf{114} (2017), no.~4, 684--732.
\href{https://doi.org/10.1112/plms.12022}{doi:10.1112/plms.12022}.

\bibitem{Kedlaya}
K. S. Kedlaya,
\emph{Analytic Number Theory}, Chapter~18, Theorem~18.2,
\url{https://kskedlaya.org/ant/chapter-18.html}.

\bibitem{KimSarnak}
H. Kim and P. Sarnak,
\emph{Refined estimates towards the Ramanujan and Selberg conjectures},
J. Amer. Math. Soc. \textbf{16} (2003), 175--181,
Appendix~2 to H. Kim, \emph{Functoriality for the exterior square of
$\mathrm{GL}_4$ and the symmetric fourth of $\mathrm{GL}_2$}.
\href{https://doi.org/10.1090/S0894-0347-02-00410-1}{doi:10.1090/S0894-0347-02-00410-1}.

\bibitem{LiLiu}
J. Li and J. Liu,
\emph{Theorem $(1+1.9)$ on the Goldbach Conjecture},
\href{https://arxiv.org/abs/2606.05224v2}{arXiv:2606.05224v2} (2026).

\bibitem{RunboLi}
R. Li,
\emph{On Chen's theorem, Goldbach's conjecture and almost prime twins II},
Math. Rep. (Bucur.) \textbf{28(78)} (2026), nos.~1--2, 39--61.
\href{https://doi.org/10.59277/mrar.2026.28.78.1.2.39}{doi:10.59277/mrar.2026.28.78.1.2.39};
\href{https://arxiv.org/abs/2405.05727v4}{arXiv:2405.05727v4}.

\bibitem{Lichtman}
J. D. Lichtman,
\emph{Primes in arithmetic progressions to large moduli, and Goldbach
beyond the square-root barrier},
\href{https://arxiv.org/abs/2309.08522v1}{arXiv:2309.08522v1} (2023).

\bibitem{MaynardI}
J. Maynard,
\emph{Primes in arithmetic progressions to large moduli I: Fixed residue
classes},
\href{https://arxiv.org/abs/2006.06572v2}{arXiv:2006.06572v2} (2021).

\bibitem{MaynardII}
J. Maynard,
\emph{Primes in arithmetic progressions to large moduli II: Well-factorable
estimates},
\href{https://arxiv.org/abs/2006.07088v1}{arXiv:2006.07088v1} (2020).

\bibitem{Pascadi}
A. Pascadi,
\emph{Large sieve inequalities for exceptional Maass forms and the
greatest prime factor of $n^2+1$}, Forum Math. Pi \textbf{14} (2026), e8.
\href{https://doi.org/10.1017/fmp.2026.10025}{doi:10.1017/fmp.2026.10025};
\href{https://arxiv.org/abs/2404.04239v3}{arXiv:2404.04239v3}.

\bibitem{Shiu}
P. Shiu,
\emph{A Brun--Titchmarsh theorem for multiplicative functions},
J. Reine Angew. Math. \textbf{313} (1980), 161--170.
\href{https://doi.org/10.1515/crll.1980.313.161}{doi:10.1515/crll.1980.313.161}.

\bibitem{Suzuki}
Y. Suzuki,
\emph{The Rosser--Iwaniec sieve}, unpublished notes, 14 March 2023,
102 pp., Proposition~11.8 and Lemma~14.7.
\url{https://www.suzuyu1729.com/notes/NagoyaSeminar2022/RosserIwaniec.pdf}.

\bibitem{TaoSW}
T. Tao,
\emph{254A, Notes 2: Complex-analytic multiplicative number theory},
9 December 2014, Exercise~64,
\url{https://terrytao.wordpress.com/2014/12/09/254a-notes-2-complex-analytic-multiplicative-number-theory/}.

\bibitem{Wu2004}
J. Wu,
\emph{Chen's double sieve, Goldbach's conjecture and the twin prime
problem}, Acta Arith. \textbf{114} (2004), 215--273.
\href{https://doi.org/10.4064/aa114-3-2}{doi:10.4064/aa114-3-2}.

\bibitem{Wu}
J. Wu,
\emph{Chen's double sieve, Goldbach's conjecture and the twin prime
problem, 2}, Acta Arith. \textbf{131} (2008), 367--387.
\href{https://arxiv.org/abs/0709.3764v1}{arXiv:0709.3764v1}.

\bibitem{Young}
M. P. Young,
\emph{Explicit calculations with Eisenstein series},
\href{https://arxiv.org/abs/1710.03624v2}{arXiv:1710.03624v2} (2017).
\end{thebibliography}
